% Generated from maps.tex by code/build.py
% A.1--A.2 and B.1--B.5.  Numbered statements are explicit for the result register.
\subsection{A.1. The integral quotient and its three-dimensional target}

We write $g,u,w,d$ for the ordered homology basis
$(\widehat\gamma,\widehat u,\widehat w,\widehat\delta)$ in this
section, and $\gamma$ for its first dual covector.  Thus $d$ is a
vector, whereas $\gamma$ is a covector.  The cusp coordinate is
$q=1/t$.  Positive target orientations are the complex orientation
on the base followed by the positive $\gamma$-circle orientation.
The complex orientation on $X$ fixes the source orientation.

\paragraph{Proposition A.1.1 (Proved).}
For the markings and the clockwise root generators in the prompt,
the homological matrices and twisting vectors are
\[
 A_1=\begin{pmatrix}1&0&0&0\\6&0&1&0\\-6&-1&-1&0\\-2&1&0&1\end{pmatrix},
 \qquad
 A_2=\begin{pmatrix}1&0&0&0\\0&0&-1&0\\-6&1&0&0\\3&0&1&1\end{pmatrix},
\]
\[
 M_0=\begin{pmatrix}1&0&0&0\\0&1&0&0\\0&1&1&0\\-1&0&0&1\end{pmatrix},
 \qquad v_1=g+2u-4w,\qquad v_2=-g-3u+3w.
\]
The common invariant vectors and covectors are $\mathbb Zd$ and
$\mathbb Z\gamma$, respectively.  The smallest invariant primitive
sublattice containing the cusp vanishing lattice is
$L=\mathbb Z\langle u,w,d\rangle=\ker\gamma$.
The formula
\[
 [z,\Pi(z)x]\longmapsto([z],\gamma(x)\bmod\mathbb Z)
 \tag{A.1.1}
\]
defines a real-analytic fibrewise group epimorphism
$J\longrightarrow B^\circ\times S^1$ with connected kernel $T_L$.
At the finite fillings it descends to a smooth proper $T^3$-bundle
over a Seifert solid torus, of multiplicities $3$ and $4$ and
translation residues $1$ and $-1$.

\begin{proof}
The first two displayed matrices are the inverse transposes of
the matrices printed on page 2 of \cite{mapsAC26}; their product
has inverse $M_0$.  Direct integral multiplication gives
$A_1^3=A_2^4=A_1A_2M_0=I$ and $A_jv_j=v_j$.
Solving the two fixed-vector equations gives $(0,0,0,d_0)$.
For a row vector $(a,b,c,e)$, the first fixed-covector equation
gives $b=2c,e=3c$ and the second gives $b=c,e=2c$;
therefore $b=c=e=0$.  Also
\[
 (M_0-I)(a,b,c,e)=(0,0,b,-a),\qquad
 A_1w=u-w.
\]
These identities prove the assertions about $L$ and the cusp
lattice $K=\mathbb Z\langle w,d\rangle$.  Their primitivity follows
from the displayed bases.  The flat real marking identifies the
deck action on a fibre with $x\mapsto A_jx$.  Since
$\gamma A_j=\gamma$, (A.1.1) descends.  Its kernel is connected
because $\gamma$ is primitive.

On a root disc the affine generator is
\[
 (s,x)\longmapsto
 (e^{-2\pi i/m}s,A_jx+v_j/m),\qquad (m,\kappa)=(3,1),(4,-1).
\]
It induces $(s,\theta)\mapsto
(e^{-2\pi i/m}s,\theta+\kappa/m)$ on the character quotient.
This action is free since $\gcd(\kappa,m)=1$.  Writing
$a=e^{2\pi i\theta}$ and choosing $r$ with $r\kappa\equiv1\pmod m$,
the invariant coordinates $s a^r$ and $a^m$ identify its quotient
with $D^2\times S^1$.  The base coordinate becomes
$(s a^r)^m(a^m)^{-r}$, and hence has multiplicity $m$ at its
central circle.  Before the finite quotient the map is a trivial
$T_L$-bundle.  The free actions on source and target are
equivariant for that map, so it descends locally to a smooth
$T_L$-bundle.  Properness follows from compactness of $T_L$.
The logarithmic translation used to glue the root chart to $J$
changes the character by its $\gamma$-component.  Applying
$\gamma$ to the actual gluing translation therefore gives the
target gluing, proving agreement on every overlap.
The exact checks, including the exhaustive mod-two fixed-vector
calculation used below, are in \texttt{code/maps\_exact.py}.
\end{proof}


Engel's primitive invariant covector is named $\psi$ in his
Lemma 3.8, labelled \texttt{covect} in the supplied source
\cite{mapsEngel}.  It is not named $\varepsilon$ there.  In his
common marking $\mathcal B_0$, the integral comparison is
\[
 P_E=\begin{pmatrix}0&-2&-1&-2\\2&-1&0&-1\\0&2&1&3\\-1&0&0&0\end{pmatrix},
 \qquad \psi P_E=-\gamma.
 \tag{A.1.2}
\]
Thus the quotient characters agree after reversing the target
circle orientation.  With the inverse root generators the
translation vectors are $a_{j-1}=-P_Ev_j/m_j$.
This is the precise sign convention in the supplied family
return, B.2.2; it explains why equality of the unnamed primitive
covectors without a sign convention would be insufficient.

\paragraph{Proposition A.1.2 (Refuted as stated; corrected).}
The literal character has no continuous extension at either
triple point or at generic points of two of the three double
curves.  It has the constant limit $0$ on the smooth locus of
$W$ and on the third double curve away from the triple points.
In the chart $q=z_0z_1z_2$, $x_1=z_1$, $x_2=z_2$, the extendible
double curve is $z_0=z_1=0$, $z_2\ne0$.  The fixed curve of the
$\mathbb C^*$-action in direction $d=e_2$ is instead
$z_0=z_2=0$, $z_1\ne0$, and is one of the nonextendible curves.
Consequently the parenthesis in A.1 identifying the extendible
curve with the fixed curve is false.

\begin{proof}
Put $q=e^{-2\pi T}$ on a positive radial ray and use the exact
cusp periods
\[
 Z_T=iTB_0+C(q),\qquad
 B_0=\begin{pmatrix}0&1\\-1&0\end{pmatrix},
\]
where $C$ is bounded and holomorphic at $0$.  For a fixed
$a\in\mathbb R^2$ the representative
$x_T(a)=\exp(2\pi iZ_Ta)$ has character $a_1\bmod\mathbb Z$ and
logarithmic position
\[
 \frac{\log|x_T(a)|}{\log|q|}
 =B_0a+T^{-1}\operatorname{Im}C(q)a.
 \tag{A.1.3}
\]
The choices $a=(-1/4,1/4)$ and $a'=(-1/2,1/4)$ have logarithmic
positions tending to $(1/4,1/4)$ and $(1/4,1/2)$.
Both are interior to the triangle with vertices $0,e_1,e_2$,
so all three chart coordinates tend to $0$.  Their character
values are $3/4$ and $1/2$.  This proves nonextension at the
first triple point.  The other triangle type has vertices
$e_1,e_2,e_1+e_2$.  Its coordinates can be taken as
$x_1x_2/q,q/x_2,q/x_1$.  Positions $(3/4,3/4)$ and
$(3/4,1/2)$ lie in its interior, give distinct values of $-y_2$
modulo one, and converge to its origin.  Translation by the
lattice preserves triangle type; these are precisely the two
triple points of $W$.

For arbitrary convergent paths in the first chart, inversion
of the exact logarithmic period matrix gives
\[
 \theta_\gamma=-\frac{\log|z_2|}{\log|q|}+o(1)\pmod{\mathbb Z}.
 \tag{A.1.4}
\]
Here the error tends to zero whenever the logarithmic ratios
stay bounded; they do stay bounded near a fixed point of a
normal-crossing stratum because $q=z_0z_1z_2$ and the three
coordinates are bounded.  If $z_0,z_1\to0$ and $z_2\to c\ne0$,
the right side tends to zero.  If $z_0,z_2\to0$ and $z_1=c\ne0$,
the paths $z_2=q^a,z_0=q^{1-a}/c$, $0<a<1$, have limits
$-a\pmod1$ at the same double point.  The same construction
with $z_0=c$ proves nonextension on $z_1=z_2=0$.
At a smooth point of a component the ratio in (A.1.4) tends
to $0$ or $1$, giving zero modulo one.  Integer translates
of these charts give the stated assertion on all of $W$.

The torus action in direction $e_2=d$ is
\[
 \lambda\cdot(z_0,z_1,z_2)
       =(\lambda^{-1}z_0,z_1,\lambda z_2).
\]
Its fixed curve is $z_0=z_2=0$.  This directly locates the
incorrect identification in the prompt.  Edge direction,
torus stabilizer, and coordinate on the quotient polygon must
not be interchanged.
\end{proof}

The nonextension above is local and pointwise.  It does not
contradict extension of the cohomology class over a cusp
neighbourhood.  The toric covering has rays $(v,1)$ for all
$v\in\mathbb Z^2$.  Meridians about its toric divisors kill
the whole cocharacter lattice: $(0,1)$ is a ray and
$(e_i,1)-(0,1)=(e_i,0)$.  The toric cover is therefore simply
connected.  Its free period deck group is $\overline\Lambda$
and hence
\[
 \pi_1(N_0)=\overline\Lambda=\Lambda/K.
 \tag{A.1.5}
\]
For the punctured cusp the map to this group kills $K$ and
the base meridian.  The latter bounds the disc given by a
point of the dense orbit on a central component.  Both
long-period covectors extend in integral $H^1$, since they
kill precisely these relations.  In contrast, an extension
of $\theta_\gamma$ to all of $X\simeq S^6$ would induce a
nonzero map $H_1(F)\to\mathbb Z$ factoring through
$H_1(X)=0$, which is impossible.  This last argument is the
global cohomological obstruction K5; it says nothing about
pointwise cusp limits.

\paragraph{Proposition A.1.3 (Proved).}
For the original data the quotient over the complement of
the cusp disc is the exterior of the $(3,4)$ torus knot,
with the specified Seifert structure.  For general admissible
$(l_0,l_1,l_2)$ the filled quotient is the lens space determined
by the following meridian slopes in one oriented integral
torus basis $(\rho_1,h)$:
\[
 \mu_1=(3,-l_1),\qquad
 \mu_2=(-4,4l_0-l_2),\qquad
 \det(\mu_1,\mu_2)=p=12l_0-4l_1-3l_2.
 \tag{A.1.6}
\]
In the lens-space convention that the second meridian is
$-q_0\mu_1+p\lambda_1$, where
$\det(\mu_1,\lambda_1)=1$, the lens parameter satisfies
$3q_0\equiv4\pmod p$.  In particular the target depends only
on the signed integer $p$, and is a three-sphere exactly
when $|p|=1$.

\begin{proof}
Let $P$ be the three-holed sphere, with oriented boundary
meridians $\rho_0+\rho_1+\rho_2=0$ and positive fibre $h$.
The two finite meridian relations are
$3\rho_1-l_1h=0$ and $4\rho_2-l_2h=0$; the cusp solid-torus
meridian is $\rho_0+l_0h$.  Eliminating $\rho_0$ gives
$\rho_2=l_0h-\rho_1$ on the annulus joining the finite solid
tori, which gives exactly (A.1.6).  Since both slopes are
primitive, this is a genus-one Heegaard splitting.  Choose
$a,b\in\mathbb Z$ with $al_1+3b=1$ and put
$\lambda_1=(a,b)$.  Direct multiplication gives
\[
 \mu_2=-q_0\mu_1+p\lambda_1,
 \quad q_0=a(4l_0-l_2)+4b,
 \quad 3q_0-ap=4.
 \tag{A.1.7}
\]
This proves the claimed lens parameter in a convention fixed
by actual meridians, including orientations.  The fundamental
group of this gluing is $\mathbb Z/|p|$.  Admissibility implies
$\gcd(p,6)=1$, so the modular inverse of $3$ used above exists.

For $(l_0,l_1,l_2)=(0,1,-1)$ the two exceptional-fibre data
are $(3,1),(4,-1)$.  They are realized by the circle action
\[
 s\cdot(z_1,z_2)=(e^{4is}z_1,e^{3is}z_2)
 \quad\text{on } S^3\subset\mathbb C^2.
\]
Its invariant map is $[z_1^3:z_2^4]$; the local circle
translation residues are $4^{-1}=1\pmod3$ and
$3^{-1}=-1\pmod4$.  Removing a regular orbit gives the
same two Seifert solid tori and their product annulus as
above.  A regular orbit is the $(3,4)$ torus knot, up to
the harmless exchange of the two coordinates.  This
identifies the manifold with its peripheral structure;
the abstract presentation $\langle x,y\mid x^3=y^{-4}\rangle$
alone would not have justified that identification.
\end{proof}

\subsection{A.2. The repaired map and the scope of its uniqueness}

\paragraph{Theorem A.2.1 (Proved).}
For every admissible member, there is a smooth map
$e:X\to Q$ over $B$ whose restriction outside an arbitrarily
small cusp neighbourhood is the literal circle quotient.
For fixed quotient gluing, its homotopy class is independent
of the cusp replacement, provided the replacement maps the
cusp neighbourhood into its designated solid torus and is
kept fixed on its outer collar.  In particular there is a
well-defined class for the original member $X\simeq S^6$.

\begin{proof}
On the cusp boundary use the product coordinates
$(q,\theta_\gamma)$ after the specified collar shear.
By (A.1.5), the boundary circle class extends to $N_0$.
Since $S^1=K(\mathbb Z,1)$, it has a continuous representative
$\theta':N_0\to S^1$.  Its restriction to the boundary is
homotopic to the given one: the two maps induce the same
homomorphism on $\pi_1$ of the boundary mapping torus.
The homotopy-extension property for a manifold with boundary
changes $\theta'$ on a collar so that it agrees exactly with
$\theta_\gamma$ there.  Relative smooth approximation gives
a smooth representative, still equal on a smaller collar.
Then $(f,\theta')$ is the required map into $D^2\times S^1$,
and agrees with the finite-filling maps on the overlap.
This proves smoothness even at the critical points of $f$:
the product map is smooth without any submersion assertion.

The inclusion $\partial N_0\to N_0$ surjects on $\pi_1$,
as (A.1.5) explicitly shows.  It follows that
$H^1(N_0;\mathbb Z)\to H^1(\partial N_0;\mathbb Z)$
is injective.  Since both spaces are connected, the long
exact sequence gives $H^1(N_0,\partial N_0;\mathbb Z)=0$.
Two circle replacements agreeing on the boundary therefore
have quotient whose relative $H^1$ class is zero, hence
have a relative homotopy.  Keeping the base coordinate
fixed gives the required homotopy of maps into the solid
torus.  Different choices of the collar size are compared
on a common smaller collar.  This proves the assertion.
\end{proof}

\paragraph{Proposition A.2.2 (Refuted as an unrestricted uniqueness assertion).}
Agreement with the literal quotient outside a neighbourhood
of $W$, without a condition on the image of that neighbourhood,
does not determine a homotopy class $S^6\to S^3$.  All twelve
classes occur among maps satisfying that bare agreement
condition.

\begin{proof}
Take a closed embedded six-ball in the interior of the cusp
neighbourhood.  Homotope a given extension on a smaller ball
to be constant there, keeping its boundary and the complement
fixed; a target coordinate ball around the image of its centre
provides the contraction.  Pinch the smaller ball to
$S^6\vee D^6$ relative to its boundary and insert any based
map $\alpha:S^6\to S^3$ on the new sphere.  The resulting
global map agrees with the original map off that ball, and
represents $[e]+\alpha$ under $[S^6,S^3]=\pi_6(S^3)$.
Distinct $\alpha$ give distinct classes.  Smooth relative
approximation preserves this conclusion.  The insertion
cannot be performed within the cusp solid torus, since that
solid torus has $\pi_6=0$; this is exactly the distinction
in Theorem A.2.1.
\end{proof}

The corrected content of \cite[Remark 8.3]{mapsEngel} is
therefore a quotient over the regular part and the two finite
fillings, followed by a replacement within its circle homotopy
class over a cusp collar.  The literal character does not
extend.  The replacement gives a canonical homotopy class
when the solid-torus image condition is retained.

